\documentclass[a4paper,10pt]{article}
\usepackage[utf8]{inputenc}
\usepackage[T1]{fontenc}
% Police vectorielle : sans elle, l'environnement peut sortir une police en image, floue à l'impression.
\usepackage{lmodern}
\usepackage[french]{babel}
\usepackage{enumitem}
% Puces pleines au premier niveau, vides au second : jamais de tiret, confondu avec un signe moins.
\frenchsetup{StandardItemLabels=true}
\setlist[itemize,1]{label=\textbullet}
\setlist[itemize,2]{label=$\circ$}
\usepackage[table]{xcolor}
\usepackage{geometry}
\usepackage{tabularx}
\usepackage{xltabular}
\usepackage{amssymb}
\usepackage{amsmath}
\usepackage{textcomp}
\usepackage{tikz}

\geometry{hmargin=0.8cm, vmargin=1.5cm}

% Commande 6 cases d'auto-évaluation sur 2 lignes
\newcommand{\caseseval}{%
    \centering
    $\square\;\square\;\square$ \\
    \vspace{2mm}
    $\square\;\square\;\square$
}
\renewcommand{\tabularxcolumn}[1]{m{#1}}

% Icônes dessinées en TikZ : aucun paquet d'icônes à installer. Décor seulement.
\newcommand{\iconeLivre}{\tikz[baseline=0pt,x=1em,y=1em]{%
    \fill[orange!80!black] (0,0.05) -- (0.46,-0.03) -- (0.46,0.72) -- (0,0.8) -- cycle;
    \fill[orange!80!black] (0.54,-0.03) -- (1,0.05) -- (1,0.8) -- (0.54,0.72) -- cycle;}}
\newcommand{\iconeQuestion}{\tikz[baseline=-0.3em]{%
    \node[circle,fill=blue!70!black,text=white,inner sep=0pt,minimum size=1em,font=\bfseries\scriptsize]{?};}}
\newcommand{\iconeCible}{\tikz[baseline=-0.3em]{%
    \draw[red!70!black,line width=0.11em] (0,0) circle (0.4em);
    \fill[red!70!black] (0,0) circle (0.15em);}}
\newcommand{\iconeAmpoule}{\tikz[baseline=0pt,x=1em,y=1em]{%
    \fill[orange!90!black] (0.3,0.52) circle (0.28);
    \fill[orange!90!black] (0.19,0.14) rectangle (0.41,0.32);
    \fill[gray!60!black] (0.2,0) rectangle (0.4,0.1);}}
\newcommand{\iconesEval}{\textcolor{red}{$\boldsymbol{\times}$}/\textcolor{green!60!black}{$\checkmark$}}

\begin{document}

\begin{center}
    \LARGE \textbf{\iconeLivre \quad FICHE MÉMO}\\[0.5em]
    \large \textbf{{{NIVEAU}} - {{CHAPITRE}} - {{TITRE}}}
\end{center}

\vspace{0.2cm}

% Colonne Conseil facultative : sans elle, |m{2cm}|X|X|, sans la dernière cellule, \multicolumn{3}.
% arraystretch : 2.2 par défaut ; 1.8 si la fiche gagne à être plus dense.
% Le tableau est continu : les objectifs s'enchainent et la coupure de page est automatique.
\renewcommand{\arraystretch}{2.2}
\begin{center}
    \begin{xltabular}{\dimexpr\textwidth-5pt\relax}{|m{2cm}|X|X|X|}
        \hline
        \rowcolor{gray!10}
        \centering \scriptsize \textbf{\iconesEval{} Auto-évaluation} &
        \centering \textbf{\iconeQuestion{} Question} &
        \centering \textbf{\iconeCible{} Réponse} &
        \centering \textbf{\iconeAmpoule{} Stratégie / Conseil} \arraybackslash \\
        \hline
        \endhead

        % --- BLOC OBJECTIF (à répéter pour chaque objectif du cours) ---
        \multicolumn{4}{|c|}{\cellcolor{blue!10}\parbox{\dimexpr\textwidth-20pt\relax}{\centering \textit{\textbf{Objectif : {{INTITULÉ_DE_L_OBJECTIF}}}}}} \\
        \hline

        % --- LIGNE (à répéter pour chaque question atomique) ---
        \caseseval & {{QUESTION_ATOMIQUE_OU_TIKZ}} & {{REPONSE_COURTE}} & \footnotesize \textit{{{CONSEIL}}} \\ \hline

    \end{xltabular}
\end{center}
\end{document}
